\documentclass[11pt,a4paper]{article}
\usepackage[utf8]{inputenc}
\usepackage[italian]{babel}
\usepackage{geometry}
\usepackage{titlesec}
\usepackage{enumitem}
\usepackage{booktabs}
\usepackage{tabularx}
\usepackage{color}
\usepackage{colortbl}
\usepackage{tcolorbox}
\usepackage{fancyhdr}
\usepackage{graphicx}
\usepackage{amssymb}

\geometry{margin=1.5cm}

\definecolor{masterpurple}{RGB}{80, 0, 120}
\definecolor{masterlightpurple}{RGB}{120, 60, 160}
\definecolor{iconianblue}{RGB}{0, 150, 180}

\titleformat{\section}
  {\normalfont\Large\bfseries\color{masterpurple}}
  {}{0em}{}[\titlerule]

\titleformat{\subsection}
  {\normalfont\large\bfseries\color{masterpurple}}
  {}{0em}{}

\pagestyle{fancy}
\fancyhf{}
\renewcommand{\headrulewidth}{1pt}
\renewcommand{\footrulewidth}{1pt}
\fancyhead[L]{Master Generale}
\fancyhead[C]{Figures of a future past}
\fancyhead[R]{Mensa Italia 2025}
\fancyfoot[C]{\thepage}

\begin{document}

\begin{center}
\Huge\textbf{\textcolor{masterpurple}{Figures of a future past}}\\[0.5cm]
\Large\textbf{Pannello Master Generale}
\end{center}

\begin{tcolorbox}[colback=masterlightpurple!20, colframe=masterpurple, title=\textbf{Premessa Narrativa}]
Nel sistema Tartarus, ai limiti della regione della Shackleton Expanse, è stata scoperta un'anomalia subspaziale che emette segnali iconiani. Gli Iconiani, leggendari per i loro portali in grado di connettere mondi distanti istantaneamente, hanno lasciato dietro di sé solo frammenti della loro avanzata tecnologia.

Questa scoperta ha riattivato un artefatto Nel sistema Tartarus, ai limiti della regione della Shackleton Expanse, è stata scoperta un'anomalia subspaziale che emette segnali iconiani. Gli Iconiani, leggendari per i loro portali in grado di connettere mondi distanti istantaneamente, hanno lasciato dietro di sé solo frammenti della loro avanzata tecnologia.

Questa scoperta ha riattivato un artefatto iconiano che minaccia di risvegliare una rete di portali in tutta la galassia. Ogni fazione coinvolta nell'avventura ha obiettivi unici e motivazioni che guidano le loro azioni, creando tensioni naturali che possono sfociare in conflitti o alleanze instabili.

L'anomalia non è solo una reliquia del passato: è un sistema iconiano semi-intelligente che manipola attivamente segnali e visioni. Questa "intelligenza incomprensibile" osserva le azioni delle fazioni e agisce di conseguenza, rendendo le intenzioni dell'artefatto un mistero da decifrare.
\end{tcolorbox}

\section{Struttura dell'Avventura}

\begin{tcolorbox}[colback=masterlightpurple!10, colframe=masterpurple]
\begin{enumerate}[leftmargin=*]
\item \textbf{Introduzione (15 minuti)}
   \begin{itemize}
   \item Ogni tavolo riceve un briefing sulla propria missione
   \item I giocatori si ambientano nei loro ruoli
   \item Si discutono le strategie iniziali
   \end{itemize}

\item \textbf{Missioni Iniziali (60 minuti)}
   \begin{itemize}
   \item Ogni tavolo affronta una sfida specifica legata alla propria fazione
   \item Ricevono frammenti di informazioni sull'anomalia
   \item Prime comunicazioni tra tavoli negli ultimi 10 minuti
   \end{itemize}

\item \textbf{Svolta (90 minuti)}
   \begin{itemize}
   \item L'anomalia si intensifica creando instabilità
   \item Un portale si attiva parzialmente, minacciando un pianeta chiave
   \item I tavoli devono collaborare per decifrare i segnali iconiani
   \end{itemize}

\item \textbf{Conclusione (75 minuti)}
   \begin{itemize}
   \item Scelta finale: stabilizzare, distruggere o attivare il portale
   \item Ogni tavolo contribuisce al risultato globale
   \item Possibili alleanze o conflitti determinano il finale
   \end{itemize}
\end{enumerate}
\end{tcolorbox}

\section{I Tavoli e le Loro Motivazioni}

\begin{tcolorbox}[colback=masterlightpurple!10, colframe=masterpurple]
\begin{description}[leftmargin=*]
\item[USS Venture (Federazione)] Esplorare e proteggere l'artefatto iconiano, prevenendo l'abuso della tecnologia. Concentrata su QDP Alpha (conoscenza).

\item[RRW Thalassar (Romulano-Ferengi)] Alleanza instabile che cerca ricchezza (Ferengi) e vantaggio strategico (Romulani). Concentrata su QDP Beta (tecnologia, potere).

\item[SS Halcyon (Nave Mineraria)] Sopravvivere agli effetti dell'anomalia e proteggere Europa Nova dalla distruzione. Buon equilibrio tra QDP Alpha e Gamma.

\item[HSS Ironclad (Nave Militare)] Neutralizzare la minaccia iconiana e garantire la sicurezza dei pianeti dell'Alleanza. Concentrata su QDP Gamma (sistemi difensivi).
\end{description}
\end{tcolorbox}

\section{Quantum Data Points (QDP)}

\begin{tcolorbox}[colback=masterlightpurple!10, colframe=masterpurple]
\textbf{Stato Iniziale:}

\begin{tabularx}{\textwidth}{|X|X|X|X|}
\hline
\rowcolor{masterlightpurple!30} \textbf{Tavolo} & \textbf{QDP Alpha} & \textbf{QDP Beta} & \textbf{QDP Gamma} \\
\hline
USS Venture & 2 & 0 & 0 \\
\hline
RRW Thalassar & 0 & 2 & 0 \\
\hline
SS Halcyon & 1 & 0 & 1 \\
\hline
HSS Ironclad & 0 & 0 & 2 \\
\hline
\end{tabularx}

\vspace{0.5cm}
\textbf{Tipi di QDP:}
\begin{description}[leftmargin=*]
\item[Alpha] Conoscenza, analisi scientifica, comprensione
\item[Beta] Tecnologia, potere, comunicazione, influenza
\item[Gamma] Risorse fisiche, materiali, sistemi difensivi
\end{description}
\end{tcolorbox}

\section{Enigmi Interconnessi}

\begin{tcolorbox}[colback=masterlightpurple!10, colframe=masterpurple]
\begin{description}[leftmargin=*]
\item[La Chiave Frammentata] Ogni nave scopre un quarto di un "codice di attivazione" iconiano. L'artefatto ha deliberatamente diviso il codice, valutando la capacità delle diverse specie di cooperare.

\item[Il Coro Galattico] L'artefatto emette quattro frequenze diverse, ciascuna rilevabile solo da una delle navi. Quando combinate, formano una "melodia" che può attivare o stabilizzare l'artefatto.

\item[Il Paradosso dei Mondi Possibili] Ogni nave sperimenta una linea temporale alternativa in cui l'artefatto è stato usato in modo diverso. L'artefatto sta valutando quale fazione ha la visione più "compatibile" con gli ideali iconiani.

\item[I Guardiani dell'Eredità] Quattro entità spettrali iconiane appaiono simultaneamente, una su ciascuna nave. Ciascuna entità rappresenta un "Guardiano" di un aspetto diverso della civiltà iconiana.
\end{description}
\end{tcolorbox}

\section{Gestione delle Manifestazioni Iconiane}

\begin{tcolorbox}[colback=iconianblue!20, colframe=iconianblue, title=\textbf{Manifestazioni Iconiane}]
\begin{description}[leftmargin=*]
\item[Sensoriali] Echi psionici, distorsioni temporali localizzate, risonanze subspaziali udibili

\item[Fisiche] Cristallizzazioni spontanee, proiezioni cartografiche, replicazione materiale

\item[Tecnologiche] Interferenze sistemiche, alterazioni energetiche, sincronizzazioni quantiche forzate

\item[Narrative] Frammenti di storia iconiana, interazioni con spettri iconiani, connessioni oniriche
\end{description}

\textbf{Utilizzo strategico:}
\begin{itemize}[leftmargin=*]
\item Fornire indizi cruciali quando i giocatori sono bloccati
\item Creare tensione e urgenza in momenti chiave
\item Facilitare la comunicazione tra tavoli in modo narrativo
\item Sviluppare il mistero degli iconiani e le loro intenzioni
\end{itemize}
\end{tcolorbox}

\section{Meccanica del Teletrasporto Internavale}

\begin{tcolorbox}[colback=masterlightpurple!10, colframe=masterpurple]
\textbf{Costo base:} 3 QDP totali

\textbf{Risoluzione:} I Master tirano 2d6 con i seguenti modificatori:
\begin{itemize}[leftmargin=*]
\item +1 per ogni QDP Alpha speso (stabilità)
\item +1 se la nave bersaglio è cooperativa
\item -1 per ogni livello di intensità dell'interferenza
\item -1 se gli scudi della nave bersaglio sono attivi
\end{itemize}

\textbf{Risultato:}
\begin{itemize}[leftmargin=*]
\item 10+: Successo completo
\item 7-9: Successo con complicazioni
\item 4-6: Problemi seri (personale a rischio)
\item 3 o meno: Fallimento catastrofico
\end{itemize}

\textbf{Capacità speciali per fazione:}
\begin{itemize}[leftmargin=*]
\item \textbf{Federazione:} "Biofiltro Avanzato" - Rileva infiltrazioni (2 QDP)
\item \textbf{Romulani:} "Infiltrazione Mascherata" - Più difficile da rilevare (3 QDP)
\item \textbf{Nave Mineraria:} "Estrazione di Emergenza" - Recupera personale in pericolo (4 QDP)
\item \textbf{Nave Militare:} "Intercettazione Teletrasporto" - Reindirizza teletrasporti nemici (3 QDP)
\end{itemize}
\end{tcolorbox}

\section{Eventi Globali e Impulsi Quantici}

\begin{tcolorbox}[colback=masterlightpurple!10, colframe=masterpurple]
\textbf{Eventi Globali:}
\begin{description}[leftmargin=*]
\item[Riattivazione Parziale Portale] Un punto dello spazio diventa instabile con rischio di anomalie temporali
\item[Collasso Subspaziale] Comunicazioni tra tavoli fortemente disturbate
\item[Messaggio Iconiano Frammentario] Ogni nave riceve un frammento di un messaggio più ampio
\item[Modifiche alle Condizioni di Scansione] Cambiano le regole di generazione QDP
\end{description}

\textbf{Impulsi Quantici:}
\begin{description}[leftmargin=*]
\item[Debole] Leggera distorsione, generando 1-2 QDP casuali per ogni nave
\item[Medio] Disturbo subspaziale, riassegnando QDP tra le tipologie
\item[Forte] Sincronizzazione quantica, trasferendo QDP tra tavoli
\item[Critico] Attivazione parziale di un portale con effetti imprevedibili
\end{description}
\end{tcolorbox}

\section{Cronometro delle Fasi}

\begin{tcolorbox}[colback=masterlightpurple!10, colframe=masterpurple]
\begin{tabularx}{\textwidth}{|l|X|l|}
\hline
\rowcolor{masterlightpurple!30} \textbf{Fase} & \textbf{Attività} & \textbf{Durata} \\
\hline
Introduzione & Briefing, presentazione personaggi, strategie iniziali & 15 minuti \\
\hline
Missioni Iniziali & Sfide fazione-specifiche, raccolta informazioni & 60 minuti \\
\hline
Svolta & Intensificazione anomalia, attivazione parziale & 90 minuti \\
\hline
Conclusione & Decisione finale sul destino dell'artefatto & 75 minuti \\
\hline
\end{tabularx}

\textbf{Promemoria temporali:}
\begin{itemize}[leftmargin=*]
\item Segnalare quando mancano 15 minuti alla fine di ogni fase
\item Negli ultimi 10 minuti della fase 2, iniziare comunicazioni tra tavoli
\item Utilizzare /final\_countdown per gestire il conto alla rovescia finale
\end{itemize}
\end{tcolorbox}

\section{Possibili Stati Finali dell'Artefatto}

\begin{tcolorbox}[colback=masterlightpurple!10, colframe=masterpurple]
\begin{description}[leftmargin=*]
\item[Stabilizzazione] (predominanza QDP Alpha) L'artefatto viene contenuto e messo in stasi. La rete di portali rimane dormiente. Previene rischi immediati ma limita la conoscenza.

\item[Attivazione Controllata] (predominanza QDP Beta) L'artefatto viene parzialmente attivato. Portali selezionati diventano operativi. Offre vantaggi tecnologici ma con rischi moderati.

\item[Estrazione] (predominanza QDP Gamma) L'artefatto viene parzialmente smantellato. Risorse preziose vengono estratte. Genera benefici immediati ma potenziali instabilità.

\item[Distruzione] (squilibrio tra QDP) L'artefatto diventa instabile e implosivo. Previene il ritorno degli Iconiani ma causa danni locali. Soluzione drastica con conseguenze potenzialmente catastrofiche.

\item[Risveglio Completo] (bilanciamento perfetto di QDP) L'artefatto si attiva completamente. La rete di portali si riattiva. Gli Iconiani potrebbero ritornare.
\end{description}
\end{tcolorbox}

\section{Comandi del Bot}

\begin{tcolorbox}[colback=masterlightpurple!10, colframe=masterpurple]
\begin{description}[leftmargin=*]
\item[/register] Registrazione come Master
\item[/status] Stato corrente dell'avventura
\item[/startadventure] Avviare l'avventura (Fase 1)
\item[/nextphase] Passare alla fase successiva
\item[/generate\_impulse] Generare un impulso quantico
\item[/global\_event] Attivare un evento globale
\item[/qdp\_status] Vedere lo stato dei QDP di tutti i tavoli
\item[/activate\_enigma <tipo>] Attivare un enigma (key, chorus, network, guardians)
\item[/manifest <numero> <tavolo>] Attivare manifestazione iconiana
\item[/final\_countdown <minuti>] Avviare il conto alla rovescia finale
\item[/upload\_media] Caricare contenuti multimediali per gli enigmi
\end{description}
\end{tcolorbox}

\section{Suggerimenti per la Narrazione}

\begin{tcolorbox}[colback=masterlightpurple!10, colframe=masterpurple]
\begin{itemize}[leftmargin=*]
\item Utilizzare le manifestazioni iconiane per fornire indizi cruciali quando i giocatori sono bloccati
\item Creare tensione attraverso countdown visibili e conseguenze chiare
\item Facilitare la comunicazione tra tavoli, specialmente se si stanno isolando
\item Alterare leggermente i messaggi tra tavoli per creare confusione strategica
\item Favorire momenti di collaborazione senza forzarli
\item Assicurarsi che ogni tavolo abbia momenti di protagonismo nell'avventura
\item Adattare le manifestazioni al tema di ciascuna nave (scientifica/diplomatica/mineraria/militare)
\item Creare un senso di coerenza attraverso le varie fasi dell'avventura
\item Celebrare i momenti di brillante cooperazione o strategia tra i tavoli
\end{itemize}
\end{tcolorbox}

\section{Tabella dei Tiri di Convergenza}

\begin{tcolorbox}[colback=masterlightpurple!10, colframe=masterpurple]
\begin{tabularx}{\textwidth}{|l|X|}
\hline
\rowcolor{masterlightpurple!30} \textbf{Risultato} & \textbf{Conseguenza} \\
\hline
≤ 4: Fallimento Catastrofico & Le azioni causano conseguenze disastrose e incontrollate \\
\hline
5-7: Fallimento Parziale & Impatto limitato e potenzialmente problematico \\
\hline
8-10: Successo Parziale & Successo con limitazioni o effetti collaterali \\
\hline
11-13: Successo & Azioni pienamente riuscite come pianificato \\
\hline
≥ 14: Successo Completo & Successo straordinario, oltre le aspettative \\
\hline
\end{tabularx}

\textbf{Modificatori per Tavolo:}
\begin{itemize}[leftmargin=*]
\item \textbf{Federazione:} +1 per ogni QDP Alpha
\item \textbf{Romulani-Ferengi:} +1 per ogni QDP Beta
\item \textbf{Nave Mineraria:} +1 per ogni QDP Gamma
\item \textbf{Nave Militare:} +1 per ogni QDP (max +8)
\end{itemize}
\end{tcolorbox}

\section{Gli Iconiani - Informazioni di Background}

\begin{tcolorbox}[colback=iconianblue!20, colframe=iconianblue]
\begin{itemize}[leftmargin=*]
\item Civiltà fiorita oltre 200.000 anni fa, soprannominati "Demoni dell'Aria e dell'Oscurità"
\item Aspetto fisico: umanoidi alti con pelle translucida iridescente e occhi completamente neri
\item Tecnologia: dispositivi semi-senzienti, manipolazione quantica della realtà
\item Portali: collegamenti diretti attraverso il tessuto subspaziale per viaggi istantanei
\item Caduta: attacco coordinato da specie soggette che temevano il loro potere
\item Possibile sopravvivenza: piccoli gruppi potrebbero essere sopravvissuti in stasi o forme evolute
\item Eredità: rovine e artefatti sparsi, varie specie che rivendicano discendenza parziale
\end{itemize}
\end{tcolorbox}

\section{Quadro delle Interazioni tra Tavoli}

\begin{tcolorbox}[colback=masterlightpurple!10, colframe=masterpurple]
\begin{tabularx}{\textwidth}{|l|X|}
\hline
\rowcolor{masterlightpurple!30} \textbf{Interazione} & \textbf{Potenziale Narrativo} \\
\hline
Federazione + Mineraria & Alleanza basata su analisi scientifica e pragmatismo \\
\hline
Federazione + Militare & Tensione tra esplorazione e sicurezza \\
\hline
Federazione + Romulani & Sospetto storico ma potenziale collaborazione scientifica \\
\hline
Romulani + Mineraria & Potenziale sfruttamento commerciale delle risorse \\
\hline
Romulani + Militare & Alta tensione e potenziale conflitto diretto \\
\hline
Mineraria + Militare & Alleanza naturale per proteggere colonie indipendenti \\
\hline
\end{tabularx}
\end{tcolorbox}

\section{Checklist per l'Avvio della Partita}

\begin{tcolorbox}[colback=masterlightpurple!10, colframe=masterpurple]
\begin{itemize}[leftmargin=*]
\item Verificare la registrazione di tutti i Master sul bot Telegram
\item Distribuire i pannelli di gioco a ciascun Master
\item Verificare la disponibilità dei materiali per gli enigmi (immagini, audio)
\item Spiegare il funzionamento dei QDP e del teletrasporto
\item Chiarire il meccanismo di interazione tra tavoli
\item Stabilire segnali non verbali per la gestione temporale
\item Ricordare a tutti i Master le capacità uniche delle loro navi
\item Effettuare un breve test di comunicazione tra tavoli
\item Sincronizzare orologi per la tempistica delle fasi
\end{itemize}
\end{tcolorbox}

\end{document}